% =====================================================================
%  §10. Метод анализа: рабочие правила с кейсами
%
%  ЖАНР (STYLE.md): учебная литература. Ни путей, ни имён программ, ни
%  номеров гипотез. Прослеживаемость — ключи \srckey{10.x} и приложение.
%
%  Содержательная основа: BRIEF_FOR_D_teaching_ladder.md §10.
% =====================================================================
\section{Метод анализа: рабочие правила с кейсами}
\label{sec:method}
\markboth{§\thesection. Метод анализа}{}

Части~I--III дали физику, аппарат оценивания и аппарат проверки достоверности. Этот заключительный
параграф~--- не новая математика, а свод из девяти рабочих правил, по которым строилось всё
предыдущее изложение. Каждое правило сформулировано абстрактно и сразу же проиллюстрировано
конкретным случаем из того же проекта, на данных которого построен весь курс~--- так, чтобы правило
запоминалось не как лозунг, а как приём с понятной ценой ошибки, если им пренебречь.

\begin{tcolorbox}[enhanced,colback=cPanel,colframe=cGrid,boxrule=0.8pt,arc=2pt,
                  title={\bfseries\color{cAccent} Пререквизиты},coltitle=cAccent,
                  attach boxed title to top left={xshift=6pt,yshift=-2pt},
                  boxed title style={colback=cPanel,boxrule=0pt,arc=0pt},top=10pt]
\textbf{Нужно знать:} весь курс, §0--§9~--- каждое правило опирается на аппарат и примеры,
введённые ранее, и явно на них ссылается.\\[0.3em]
\textbf{Знать НЕ нужно:} ничего нового~--- этот параграф собирает уже известное в рабочий
чек-лист.
\end{tcolorbox}

% ---------------------------------------------------------------------
\subsection{1. Модельно-независимое прежде модельного}
\label{subsec:rule1}

Прежде чем доверять выводу, полученному подгонкой сложной модели, стоит спросить: можно ли
установить тот же факт способом, не зависящим от правоты модели вообще,~--- симметрией,
периодичностью, законом сохранения энергии?

\begin{example}[: расхождение объяснено ещё до всякой подгонки]\srckey{10.1}
На одном из образцов повторные измерения, снятые в разные дни, расходились в области глубокого
подавления почти в полтора раза. Причина была установлена \emph{без единого фита}: угол $-95^\circ$
физически задаёт ту же ориентацию проводов, что и $+85^\circ$ (поляризатор не различает
направление вращения на $180^\circ$~--- §\ref{sec:electrodynamics}), а измерения в двух
конфигурациях эту симметрию не учитывали одинаково. Модельная подгонка, проведённая позже,
только подтвердила уже установленный факт.
\end{example}

% ---------------------------------------------------------------------
\subsection{2. Отрицательный контроль обязателен}
\label{subsec:rule2}

Эффект, добавленный в модель, стоит проверять не только там, где он предположительно есть, но и
там, где, по своей же физике, его быть не должно. Если он «находится» одинаково хорошо
и в контрольных условиях, подозрение падает на метод, а не на физику.

\begin{example}[: просачивание проходит контроль честно]\srckey{10.2}
Явное просачивание нужно плотной решётке (уверенное улучшение $\AIC$, §6) и \emph{не нужно}
разреженному образцу: там амплитуда просачивания стремится к нулю, а добавление параметра лишь
штрафует $\AIC$ без всякого выигрыша. Это ровно то поведение, которого требует отрицательный
контроль,~--- эффект, «находимый» одинаково везде независимо от физических условий, как правило,
оказывается артефактом метода, а не реальностью.
\end{example}

% ---------------------------------------------------------------------
\subsection{3. Предсказание фиксируется до расчёта}
\label{subsec:rule3}

Самая убедительная проверка гипотезы~--- не подгонка задним числом, а численное предсказание,
сделанное \emph{на бумаге}, до того как получен результат расчёта или измерения.

\begin{example}[: теорема, а не апостериорное объяснение]\srckey{10.3}
Теорема о выровненном анализаторе (§1) предсказала конкретное число~--- изменение информационного
критерия, близкое к нулю,~--- \emph{до} того, как модель со вторым реальным элементом тракта была
запущена. Измерено $\dAIC=+0{,}15$ (§1, §6): предсказание подтвердилось количественно, а не
только по знаку. Такая проверка гораздо строже, чем подбор объяснения под уже известный результат.
\end{example}

% ---------------------------------------------------------------------
\subsection{4. Триангуляция}
\label{subsec:rule4}

Один путь к результату~--- гипотеза. Два и более независимых пути, сходящихся к одному и тому же
числу,~--- уже результат, которому можно доверять существенно больше.

\begin{example}[: просачивание подтверждено тремя разными способами]\srckey{10.4}
Существование и величина просачивания подтверждены: (а) явным включением его в подгонку
(§6); (б) энергетическим балансом развёртки без всякой подгонки спектра, по формуле~$\cos^4\theta$
из §0; (в) прямым наблюдением временного сдвига пика импульса на осциллограмме (§0,
$\tleak\approx0{,}19$~пс, ключ~0.1). Каждый путь использует разные данные и разные предположения;
их совпадение~--- физический аргумент, а не совпадение арифметики.
\end{example}

% ---------------------------------------------------------------------
\subsection{5. Различать «эффект» и «параметр»}
\label{subsec:rule5}

Это правило курс формулировал по частям в §3, §6, §8 и §9~--- здесь оно собрано явно, как
самостоятельный пункт метода.

\begin{tcolorbox}[enhanced,colback=cS1!7,colframe=cS1,boxrule=0pt,leftrule=3pt,arc=1pt]
«Модель с этим членом лучше описывает данные» и «численное значение этого параметра надёжно
измерено»~--- разные утверждения. Курс видел это различение на показателе рассеяния (эффект
реален, значение плывёт в разы, §3), на перекрёстном члене (улучшает подгонку, не проходит
проверку однородности, §6), на амплитуде просачивания (хорошо описывает обучающие данные, не
экстраполируется в held-out проверке, §8) и на эффективном диаметре (нужен данным, но
неотделим от просачивания без независимой фиксации, §9). Каждый раз первое утверждение было
верным, а второе~--- нет, и различить их можно было только специальной проверкой, а не более
внимательным взглядом на тот же результат.
\end{tcolorbox}

% ---------------------------------------------------------------------
\subsection{6. При вырождении не подгонять оба параметра}
\label{subsec:rule6}

Если корреляция двух параметров близка к предельной (§9), их совместная свободная подгонка не
даёт двух чисел~--- она даёт одно число (комбинацию) и иллюзию второго. Правильное решение~---
зафиксировать тот из пары, для которого есть независимое обоснование значения, и подгонять только
второй, отдавая себе отчёт, что он измеряет именно комбинацию, а не чистую физическую величину
по отдельности (§9, §6).

% ---------------------------------------------------------------------
\subsection{7. Дешёвая диагностика перед дорогой подгонкой}
\label{subsec:rule7}

Прежде чем запускать многостартовую нелинейную оптимизацию (§7)~--- дорогую и по времени, и по
риску не туда сойтись,~--- стоит проверить, что даёт линейный гармонический разбор (§4): решение
одной линейной системы, аналитические ошибки, единственный минимум. Он не заменяет полную модель,
но даёт офсет угла, амплитуду и относительную фазу просачивания почти бесплатно~--- готовый
источник тёплого старта (§7) и независимую опорную точку для проверки результата дорогой
подгонки.

% ---------------------------------------------------------------------
\subsection{8. Отзывать вывод при новой информации — в тот же день и везде}
\label{subsec:rule8}

Наука отличается от догмы не отсутствием ошибок, а тем, что делают, когда ошибка обнаружена.
Правило: как только появляется информация, меняющая вывод, вывод отзывается сразу и во всех
местах, где он был опубликован,~--- а не тихо забывается или оставляется висеть рядом с новыми
данными, ему противоречащими.

\begin{example}[: отзыв вывода в тот же день, когда обнаружилась его причина]\srckey{10.5}
Курс уже описывал этот случай в §\ref{sec:electrodynamics} (§0, предупреждение об оси вращения):
вывод о связи дефицита эффективного диаметра с качеством решётки был опубликован по данным одной
из серий измерений, а вечером того же дня выяснилось, что пучок в этой серии был намеренно смещён
от оси вращения~--- из-за чего угол поворота и участок образца оказались неразделимы, и наблюдаемая
могла с равным успехом объясняться пространственной неоднородностью, а не углом. Вывод был отозван
в тот же день, во всех документах, где он был записан, с явной пометкой «не проверено», а не
исправлен молча. Это штатная операция метода, а не признание катастрофы: историю не переписывают,
рядом с прежней записью ставится поправка.
\end{example}

% ---------------------------------------------------------------------
\subsection{9. Артефакт — или не было}
\label{subsec:rule9}

Каждое число в курсе имеет источник~--- конкретный протокол измерения или расчёт, к которому
можно вернуться и проверить. Именно поэтому у курса есть служебное приложение
«Происхождение числовых примеров»: не потому, что читателю это интересно само по себе, а потому,
что дисциплина «число без источника не существует» задним числом ловит куда больше ошибок, чем
любая проверка по памяти. Ключи вида \srckey{10.x}, встречавшиеся на каждой странице,~---
не формальность, а прямое применение этого правила к самому курсу.

% ---------------------------------------------------------------------
\subsection{Свод правил}

\begin{center}
\footnotesize
\setlength{\tabcolsep}{6pt}
\begin{tabular}{@{}p{0.05\linewidth}p{0.40\linewidth}p{0.49\linewidth}@{}}
\toprule
\textbf{№} & \textbf{Правило} & \textbf{Где в курсе} \\
\midrule
1 & Модельно-независимое прежде модельного & §\ref{subsec:rule1} \\
2 & Отрицательный контроль обязателен & §\ref{subsec:rule2} \\
3 & Предсказание фиксируется до расчёта & §\ref{subsec:rule3}, теорема §1 \\
4 & Триангуляция & §\ref{subsec:rule4} \\
5 & Различать «эффект» и «параметр» & §3, §6, §8, §9 \\
6 & При вырождении не подгонять оба параметра & §5, §9 \\
7 & Дешёвая диагностика перед дорогой подгонкой & §4, §7 \\
8 & Отзывать вывод сразу и везде & §\ref{subsec:rule8}, §0 \\
9 & Артефакт — или не было & приложение «Происхождение числовых примеров» \\
\bottomrule
\end{tabular}
\end{center}

\begin{intuition}[: что объединяет все девять правил]
Ни одно из них не является специфичным для терагерцовой спектроскопии или для поляризационной
оптики~--- это правила работы с любой моделью, подгоняемой к зашумлённым данным в условиях, когда
число параметров сопоставимо с числом независимых чисел, которые содержит измерение. Курс шёл от
электродинамики (§0) к статистике оценивания (§4--§9) именно затем, чтобы эти правила не звучали
абстрактно, а имели за собой конкретную математику и конкретную цену ошибки, если ими пренебречь.
\end{intuition}

\clearpage
